% TODO make hw stylesheet
\documentclass[10pt]{article}
\usepackage[letterpaper,top=1in]{geometry}
\usepackage[dvipsnames,svgnames]{xcolor}
\usepackage[shortlabels]{enumitem}
\usepackage[framemethod=TikZ]{mdframed}
\usepackage{amsmath,amssymb,amsthm}
\usepackage[colorlinks]{hyperref}
\usepackage{multicol}
\usepackage{tabularx}

\hypersetup{
	colorlinks=true,
	linkcolor=blue,
	filecolor=magenta,
	urlcolor=magenta,
	pdftitle={MAT 534 HW}
}

\usepackage{mathtools}
\usepackage{thmtools}
\usepackage[textsize=footnotesize]{todonotes}
\usepackage{titling}

\reversemarginpar
\mdfdefinestyle{mdbluebox}{roundcorner=10pt,innerbottommargin=9pt,
	linecolor=blue,backgroundcolor=TealBlue!5,}
\declaretheoremstyle[headfont=\sffamily\bfseries\color{MidnightBlue},
mdframed={style=mdbluebox},]{thmbluebox}

\mdfdefinestyle{mdbloobox}{frametitlefont=\bfseries,innerbottommargin=8pt,
	nobreak=true,backgroundcolor=TealBlue!5,linecolor=blue,}
\declaretheoremstyle[headfont=\bfseries\color{MidnightBlue},
mdframed={style=mdbloobox},headpunct={\\[3pt]},postheadspace=0pt,]{thmbloobox}

\mdfdefinestyle{mdredbox}{frametitlefont=\bfseries,innerbottommargin=8pt,
	nobreak=true,backgroundcolor=Salmon!5,linecolor=RawSienna,}
\declaretheoremstyle[headfont=\bfseries\color{RawSienna},
mdframed={style=mdredbox},headpunct={\\[3pt]},postheadspace=0pt,]{thmredbox}

\mdfdefinestyle{mdgreenbox}{linecolor=ForestGreen,backgroundcolor=ForestGreen!5,
	linewidth=2pt,rightline=false,leftline=true,topline=false,bottomline=false,}
\declaretheoremstyle[headfont=\bfseries\sffamily\color{ForestGreen!70!black},
mdframed={style=mdgreenbox},headpunct={ --- },]{thmgreenbox}

\mdfdefinestyle{mdorangebox}{linecolor=RedOrange,backgroundcolor=RedOrange!5,
	linewidth=2pt,rightline=false,leftline=true,topline=false,bottomline=false,}
\declaretheoremstyle[headfont=\bfseries\sffamily\color{RedOrange!70!black},
mdframed={style=mdorangebox},headpunct={ --- },]{thmorangebox}

\mdfdefinestyle{mdblackbox}{linecolor=black,backgroundcolor=RedViolet!5!gray!5,
	linewidth=3pt,rightline=false,leftline=true,topline=false,bottomline=false,}
\declaretheoremstyle[mdframed={style=mdblackbox}]{thmblackbox}

\declaretheoremstyle[spaceabove=6pt,spacebelow=6pt]{basehead}

\declaretheorem[style=basehead,name=Problem]{problem}
\declaretheorem[style=thmredbox,name=Problem,numbered=no]{problem*}
\declaretheorem[style=thmbloobox,name=Setup]{setup} % for config geo
\declaretheorem[style=thmredbox,name=Example,sibling=problem]{example}
\declaretheorem[style=thmredbox,name=$\bigstar$ Problem,sibling=problem]{niceprob}
\declaretheorem[style=thmbluebox,name=Theorem,sibling=problem]{theorem}
\declaretheorem[style=thmbluebox,name=Corollary,sibling=theorem]{corollary}
\declaretheorem[style=thmbluebox,name=Proposition,sibling=theorem]{prop}
\declaretheorem[style=thmbluebox,name=Lemma,numberwithin=problem]{lemma}
\declaretheorem[style=thmbluebox,name=Theorem,numbered=no]{theorem*}
\declaretheorem[style=thmbluebox,name=Lemma,numbered=no]{lemma*}
\declaretheorem[style=thmgreenbox,name=Claim,numberwithin=problem]{claim}
\declaretheorem[style=thmgreenbox,name=Claim,numbered=no]{claim*}
\declaretheorem[style=thmblackbox,name=Remark,numberwithin=problem]{remark}
\declaretheorem[style=thmblackbox,name=Remark,numbered=no]{remark*}
\declaretheorem[style=thmgreenbox,name=Definition,sibling=theorem]{definition}
\declaretheorem[style=thmgreenbox,name=Definition,numbered=no]{definition*}

\providecommand{\ol}{\overline}
\providecommand{\quiv}{\equiv}
\providecommand{\ceil}[1]{\left\lceil #1 \right\rceil}
\providecommand{\floor}[1]{\left\lfloor #1 \right\rfloor}
\newcommand{\dd}{\mathrm{d}}
\providecommand{\ul}{\underline}
\providecommand{\eps}{\varepsilon}
\providecommand{\half}{\frac{1}{2}}
\providecommand{\CC}{\mathbb C}
\providecommand{\FF}{\mathbb F}
\providecommand{\NN}{\mathbb N}
\providecommand{\QQ}{\mathbb Q}
\providecommand{\RR}{\mathbb R}
\providecommand{\ZZ}{\mathbb Z}
\providecommand{\dg}{^\circ}
\providecommand{\ii}{\item}
\providecommand{\setdiff}{\smallsetminus}
\providecommand{\alert}[1]{{\sffamily\textbf{\textcolor{blue}{#1}}}}
\providecommand{\inv}{^{-1}}
\providecommand{\mb}[1]{\mathbf{#1}}

\newenvironment{soln}{\begin{proof}[Solution]}{\end{proof}}

\providecommand{\pow}{\operatorname{Pow}}
\providecommand{\vocab}[1]{{\textbf{\textcolor{ForestGreen}{#1}}}}
\providecommand{\lcm}{\operatorname{lcm}}
\providecommand{\abs}[1]{\left|#1\right|}

\usepackage{mathrsfs}

% place title at top right
\setlength{\droptitle}{-8ex}
\pretitle{\begin{flushright}\Large\bfseries}
\posttitle{\par\end{flushright}}
\preauthor{\begin{flushright}}
\postauthor{\end{flushright}}
\predate{\begin{flushright}}
\postdate{\end{flushright}}

\title{MAT 534 HW02}
\author{\large Aditya Pahuja}
\date{\large Due 09/17}
\begin{document}
\maketitle

\begin{problem}
    Let $H$ and $K$ be subgroups of a finite group $G$.
    Let $HK = \{hk : h\in H,\; k\in K\}$, which need not be a subgroup of $G$.
    Show that
    \begin{equation*}
	\abs{HK} = \frac{\abs H\cdot\abs K}{\abs{H\cap K}}.
    \end{equation*}
\end{problem}

\begin{soln}
    Let $h_1$ and $h_2$ be in $H$. Then, the left cosets $h_1K$ and $h_2K$
    are equal if and only if $h_2\inv h_1K = K$, which is true if and only if
    $h_2\inv h_1\in K$, hence $h_2\inv h_1\in H\cap K$.
    Defining a relation $\sim$ on $H$ by decreeing $h_1\sim h_2$ if $h_2\inv h_1$
    tells us that the equivalence classes correspond to cosets of
    the subgroup $H\cap K$ in $H$; in particular there are
    $\abs H/\abs{H\cap K}$ cosets of $H\cap K$.
    On the other hand each coset of $H\cap K$ corresponds to a coset of $K$
    by an element of $H$, according to the first sentence, and these cosets
    of course partition $HK$, so $\abs{HK} = \abs{K}\cdot (\abs H/\abs{H\cap K})$
    as desired.
\end{soln}

\begin{problem}
    Let $G$ be a group. Define a new operation on $G$ by
    \begin{equation*}
        x*y = yx.
    \end{equation*}
    \begin{enumerate}[(a)]
        \ii Show that $G$ is a group under the new operation.
	We write $G^{\text{op}}$ for the set $G$ with the new operation $*$.
	\ii Show that $G^{\text{op}}$ is isomorphic to $G$.
	\ii Suppose that $G$ acts on $X$ by a left action,
	and let $\alpha\colon G\times X\to X$ be the map which defines the action;
	in other words, if $gx$ denotes the action of $g$ on $x$, then
	$\alpha(g,x) = gx$. Show that $\alpha$ defines a
	\emph{right} action of $G^{\text{op}}$ on $X$.
    \end{enumerate}
\end{problem}

\begin{soln}
    (a) The operation is associative because
    \begin{equation*}
	(x*y)*z = (yx)*z = zyx = x*(zy) = x*(y*z).
    \end{equation*}
    Its identity is the same as that of $G$ because
    $1*x = x\cdot1 = x = 1\cdot x = x*1$,
    and the inverse of $g\in G$ is still $g\inv$, its inverse under $\cdot$,
    because $g*g\inv = g\inv\cdot g = 1 = g\cdot g\inv = g\inv*g$.\\

    (b) Consider the map $\varphi\colon G\to G^{\text{op}}$ by
    $\varphi(g) = g\inv$. This is of course a bijection;
    it is a homomorphism because
    \begin{equation*}
        \varphi(xy) = (xy)\inv = y\inv x\inv = \varphi(x)*\varphi(y).
    \end{equation*}

    (c) To check that $\alpha$ defines a right action, we need to check
    that $\alpha(1,x) = x$ for all $x$ (which is immediate because
    $\alpha$ a priori defines a left action of $G$ on $X$) and
    \begin{equation*}
        \alpha(h,\alpha(g, x)) = \alpha(g*h, x).
    \end{equation*}
    This is easy to see because we know from $\alpha$
    defining a left action of $G$ on $X$ that
    \begin{equation*}
        \alpha(h,\alpha(g,x)) = \alpha(hg,x) = \alpha(g*h, x).\qedhere
    \end{equation*}
\end{soln}

\pagebreak

\begin{problem}
    Let $G$ and $H$ be groups, and let $\Phi$ be
    the set of all functions $\phi\colon G\to H$
    with the property that $\phi(1) = 1$.
    \begin{enumerate}[(a)]
        \ii Let $\alpha\colon G\times\Phi\to\Phi$ be defined as follows:
	$\alpha(g,\phi)$ is the function that takes $x\in G$ to
	$\phi\inv(g)\phi(gx)\in H$. Prove that this defines an action
	of $G$ on $\Phi$. Is it a right action or a left action?
	\ii Describe the set of fixed points $\Phi_0$ for this action.
    \end{enumerate}
\end{problem}

\begin{soln}
    (a) $\alpha(h, \alpha(g,\phi))$ is the function sending
    $x$ to $\psi\inv(h)\psi(hx)$
    where $\psi(x) = \phi\inv(g)\phi(gx)$ so it is the function
    \begin{equation*}
        x\mapsto (\phi\inv(g)\phi(gh))\inv(\phi\inv(g)\phi(ghx)) =
	\phi\inv(gh)\phi(ghx) = \alpha(gh,x).
    \end{equation*}
    Combined with the observation $\alpha(1,\phi)$ being the map
    $x\mapsto\phi(1)\inv\phi(x) = \phi(x)$, which is $\phi$,
    we see that $\alpha$ defines a left action on $\Phi$.

    (b) A function $\phi\in\Phi$ is a fixed point for the action if
    \begin{equation*}
        \phi(x) = \alpha(g,\phi)(x) = \phi(g)\inv\phi(gx)
    \end{equation*}
    for all $x,g\in G$. By taking $x = g\inv$ we get $\phi(g\inv) = \phi(g)\inv$
    so $\phi(x) = \phi(g\inv)\phi(gx)$ for all $x,g\in G$.
    We can replace $g$ with $g\inv$ and plug in $x = gh$ for any $h\in H$
    to see that $\phi$ satisfies
    \begin{equation*}
        \phi(gh) = \phi(g)\phi(h)
    \end{equation*}
    for all $g,h\in H$; that is, $\phi$ must be a homomorphism $G\to H$
    if it is a fixed point of the action defined by $\alpha$.
    Conversely, it is easy to see that all homomorphisms are fixed points
    of the action because if $\phi$ is a homomorhpism then
    \begin{equation*}
        \alpha(g,\phi)(x) = \phi(g)\inv\phi(gx) = \phi(g\inv gx) = \phi(x)
    \end{equation*}
    for all $x\in G$. Therefore the fixed points of the action given by $\alpha$
    are precisely the group homomorphisms $\phi\colon G\to H$.
\end{soln}

\begin{problem}
    Let $p$ be a prime, and let $G$ be a $p$-group.
    Suppose that $G$ acts on a finite set $X$. Let $X_0\subseteq X$
    be the set of fixed points of $X$ under this action. Show that
    \begin{equation*}
	\abs{X_0}\quiv\abs X\bmod p.
    \end{equation*}
\end{problem}

\begin{soln}
    Let $X_0$ be the set of fixed points of $X$
    and $X_1,X_2,\dots,X_r\subseteq X$ the orbits of $X$
    with more than one element. By the orbit-stabilizer theorem,
    each $\abs{X_j}$ for $1\le j\le r$ is a divisor of $\abs G$
    (because elements of the orbit are in bijection with cosets of
    the stabilizer in $G$ of an element of the orbit)
    that is bigger than $1$, hence is power of $p$ bigger than $1$,
    so in particular is a multiple of $p$.
    The orbits also partition $X$, so
    \begin{equation*}
	\abs X = \abs{X_0} + \underbrace{\abs{X_1} + \cdots + \abs{X_r}}_{0\bmod p}
	\quiv \abs{X_0}\bmod p.\qedhere
    \end{equation*}
\end{soln}

\pagebreak

\begin{problem}
    Let $G$ be a finite group of order $n$, let $p$ be a prime dividing $n$,
    and let $\ZZ/p\ZZ$ be the cyclic group of order $p$.
    Let $\Phi$ be the set of all functions $\phi\colon\ZZ/p\ZZ\to G$
    with the property that $\phi(1) = 1$.
    \begin{enumerate}[(a)]
	\ii Show that $\abs{\Phi} = n^{p-1}$.
	\ii The group $\ZZ/p\ZZ$ acts on the set $\Phi$, as in Problem 3.
	Let $\Phi_0$ be the set of fixed points in $\Phi$
	for the action of $\ZZ/p\ZZ$. Show that the cardinality of $\Phi_0$
	is divisible by $p$.
	\ii Show that $\Phi_0$ is nonempty, and hence contains
	at least $p$ elements. Use this to prove \emph{Cauchy's theorem}:
	if $p$ is a prime that divides the order of a finite group $G$,
	then $G$ contains a subgroup isomorphic to the cyclic group $\ZZ/p\ZZ$.
    \end{enumerate}
\end{problem}

\begin{soln}
    (a) We know that $\phi(1) = 1$; the other $p-1$ elements of $\ZZ/p\ZZ$
    each have $\abs G = n$ choices for their image in $G$,
    which nets a total of $n^{p-1}$ choices.\\

    (b) The cardinality of the set of fixed points of the action is congruent
    to the cardinality of $\Phi$ mod $p$ by Problem 4,
    since $\ZZ/p\ZZ$ is a $p$-group, and we know that
    $p\mid n^{p-1} = \abs\Phi$.
    \\

    (c) As shown in Problem 3, $\Phi_0$ consists of the homomorphisms
    $\ZZ/p\ZZ\to G$. Since $\Phi_0$ contains the trivial homomorphism
    $\phi(n)\quiv 1_G$ for all $n\in\ZZ/p\ZZ$, it must contain
    at least $p-1$ other homomorphisms, none of which are trivial.
    In particular each one has trivial kernel because the only two options
    for the kernel of a homomorphism coming out of $\ZZ/p\ZZ$
    are the trivial subgroup and $\ZZ/p\ZZ$ itself, so
    each nontrivial homomorphism describes
    a way to embed $\ZZ/p\ZZ$ as a subgroup of $G$,
    proving that $G$ has a subgroup isomorphic to $\ZZ/p\ZZ$.
\end{soln}

\begin{problem}
    Let $G$ be a group, $H$ a subgroup, and let $X$ be the set of
    left cosets of $H$. $G$ acts on $X$ by the rule
    $(g,aH)\mapsto gaH$, a left action.
    \begin{enumerate}[(a)]
        \ii Show that $G$ acts transitively on $X$.
	\ii What is the stabilizer of the coset $aH\in X$ in $G$?
	\ii Show that the homomorphism $G\to S_X$
	corresponding to this action has kernel
	\begin{equation*}
	    \bigcap_{g\in G}gHg\inv.
	\end{equation*}
	\ii Suppose that $G$ is a finite group, and that
	$H$ is a proper subgroup. Suppose that $(G:H)!$ is
	not divisible by $|G|$. Show that $G$ is not simple.
	\ii Suppose that $G$ is finite, and that $(G:H)$ is
	the smallest prime dividing $\abs G$.
	Show that $H$ is normal in $G$.
    \end{enumerate}
\end{problem}

\begin{soln}
    (a) $G$ acts transitively on $X$ because
    given cosets $aH$ and $bH$ for $a,b\in G$, we have
    \begin{equation*}
	(ba\inv)aH = bH,
    \end{equation*}
    implying that the two cosets lie in the same orbit.
    \\

    (b) The stabilizer of $aH$ is the subgroup $aHa\inv$.
    It is clear that these elements all fix $aH$;
    conversely, if $gaH = aH$, then for each $h\in H$
    there is $h'\in H$ such that
    \begin{equation*}
        gah = ah' \iff g = ah'h\inv a\in aHa\inv,
    \end{equation*}
    which gives the opposite containment of the stabilizer within $aHa\inv$.
    \\

    (c) The kernel of this homomorphism is just the elements of $G$
    whose action on $X$ fixes every coset, i.e. $g$ is in the stabilizer
    $aHa\inv$ of $aH$ for all $a$ so $g\in \bigcap_{a\in G}aHa\inv$.
    \\

    (d) We prove the contrapositive.
    Suppose $G$ is simple. Then $G$ has no nontrivial normal subgroup,
    so the homomorphism $G\to S_X$ from part (c) has kernel $1$ or $G$.
    The action of $G$ is transitive by (a),
    so the homomorphism is nontrival, hence must have kernel $1$.
    The image of the homomprhism in $S_X$ is then isomorphic to $G$,
    so $\abs G$ divides $\abs{S_X} = \abs X!$, but
    $\abs X$ is the number of cosets of $H$ in $G$, which is precisely $(G:H)$,
    which means we have proven $\abs G$ divides $(G:H)!$.
    \\

    (e) Let $X$ be the set of cosets of $H$.
    If $H$ is the trivial subgroup, then it is trivially normal in $G$.
    Otherwise, $\abs G$ has at least one prime factor besides $(G:H)$,
    which must be at least $(G:H)$. Since $(G:H)!/(G:H) = ((G:H)-1)!$
    does not have any prime factors greater than or equal to $(G:H)$,
    $\abs G$ cannot divide $(G:H)!$.

    This means that the order of
    the image of the homomorphism $G\to S_X$ (as in (c)),
    being a factor of $\gcd(\abs{S_X}, \abs G) = (G:H)$,
    is either $1$ or $(G:H)$, due to the primality of $(G:H)$.
    Again, the action of $G$ on $X$ is transitive,
    so the homomorphism is nontrivial, meaning
    the image cannot have order $1$, hence must have order $(G:H)$;
    therefore the kernel has order $\abs G/(G:H) = \abs H$.
    By (c) the kernel is also a subgroup of $H$,
    so this proves that the kernel is exactly $H$,
    which is enough to show that $H$ is normal in $G$.
\end{soln}

\begin{problem}
    Let $G = GL_2(\CC)$.
    \begin{enumerate}[(a)]
        \ii Give a complete set of representatives for
	the conjugacy classes of $G$.
	\ii Given $A\in G$, how do you determine
	which element of your list $A$ is conjugate to?
    \end{enumerate}
\end{problem}

\begin{soln}
    (a) Each square matrix with complex entries
    is similar (conjugate) to a matrix in Jordan canonical form, so
    the conjugacy classes are represented by the matrices
    in Jordan canonical form that are invertible:
    \begin{equation*}
	\left\{ \begin{pmatrix}
		a&0\\0&b
	\end{pmatrix} : a,b\in\CC\setminus\{0\}, a\succ b\right\}\cup
	\left\{ \begin{pmatrix}
		a&1\\0&a
	\end{pmatrix} : a\in\CC\setminus\{0\}\right\}.
    \end{equation*}
    where $a\succ b$ means that the real part of $a$ is greater than that of $b$
    or, if the real parts are equal, then the imaginary part of $a$
    is at least as large as that of $b$ (i.e. order $a$, $b$
    as pairs of real numbers in lexicographical order)
    so that we don't double count matrices with the same multiset of
    diagonal entries (which are conjugate).

    (b) Let $A\in G$. If $A$ has two distinct eigenvalues $\lambda_1$, $\lambda_2$
    with $\lambda_1\succ\lambda_2$, then its conjugacy class is represented by
    \begin{equation*}
        \begin{pmatrix}
	    \lambda_1&0\\0&\lambda_2
        \end{pmatrix}.
    \end{equation*}
    If $\lambda_1=\lambda_2=\lambda$, then either $A = \lambda I$
    where $I$ is the identity or it is conjugate to
    \begin{equation*}
        \begin{pmatrix}
	    \lambda&1\\0&\lambda
        \end{pmatrix};
    \end{equation*}
    it can't be conjugate to $\lambda I$ since the conjugacy class
    of $\lambda I$ is just $\{\lambda I\}$ because $\lambda I\in Z(G)$.
\end{soln}










\end{document}
